\documentclass[10pt]{article}
\usepackage[T1]{fontenc}
\usepackage[utf8]{inputenc}
\usepackage[french]{babel}
\usepackage{amsmath,amssymb,mathrsfs,mathpazo}
\usepackage{geometry}
\geometry{a4paper,margin=2cm}
\usepackage{graphicx,array,tabularx,tabularray,colortbl,multirow,multicol,booktabs,enumitem,pifont,ulem,fancybox,diagbox,cancel}
\usepackage{pstricks,pstricks-add,pst-plot,pst-tree,pst-node,pst-text,pst-func,pst-math,pst-bspline,pst-3dplot,pst-eucl,pst-solides3d}
\let\pscalculate\undefined
\usepackage{pst-bezier}
\usepackage{tikz,pgfplots,tkz-tab}
\usetikzlibrary{arrows,arrows.meta,patterns,calc,positioning,trees,decorations.markings,shapes.geometric,shapes.misc,shapes.arrows,intersections,decorations.pathmorphing,decorations.pathreplacing,backgrounds,shadows,babel}
\pgfplotsset{compat=1.18}
\usepgfplotslibrary{fillbetween}
\usepackage[np]{numprint}
\usepackage{listings,makecell,pas-tableur,twoopt}
\let\barre\undefined
\usepackage[pstricks]{tablvar}
\definecolor{darkgreen}{RGB}{0,100,0}
\providecommand{\red}{\color{red}}
\providecommand{\blue}{\color{blue}}
\providecommand{\green}{\color{green}}
\providecommand{\black}{\color{black}}
\providecommand{\white}{\color{white}}
\providecommand{\gray}{\color{gray}}
\providecommand{\cyan}{\color{cyan}}
\providecommand{\magenta}{\color{magenta}}
\providecommand{\yellow}{\color{yellow}}
\providecommand{\violet}{\color{violet}}
\providecommand{\orange}{\color{orange}}
\providecommand{\euro}{\texteuro}
\providecommand{\EUR}{\texteuro}
\providecommand{\ieme}{\textsuperscript{e}}
\providecommand{\degres}{\ensuremath{^\circ}}
\providecommand{\degree}{\ensuremath{^\circ}}
\providecommand{\vv}[1]{\overrightarrow{#1}}
\providecommand{\Coord}[1]{#1}
\providecommand{\decosix}{}
\providecommand{\starredbullet}{\ensuremath{\bullet}}
\pagestyle{empty}
\setlength{\parindent}{0pt}
\usepackage{adjustbox}\DeclareUnicodeCharacter{2236}{:}\DeclareUnicodeCharacter{2190}{\ensuremath{\leftarrow}}
\providecommand{\R}{\mathbb{R}}
\providecommand{\N}{\mathbb{N}}
\providecommand{\Z}{\mathbb{Z}}
\providecommand{\Q}{\mathbb{Q}}
\providecommand{\D}{\mathbb{D}}
\providecommand{\C}{\mathbb{C}}
\providecommand{\e}{{\mathrm{e}}}
\providecommand{\vect}[1]{\overrightarrow{#1}}
\providecommand{\vectt}[1]{\overrightarrow{\mathrm{#1}}}
\providecommand{\vvt}[1]{\overrightarrow{\mathrm{#1}}}
\providecommand{\np}[2][]{{#2~#1}}
\providecommand{\vv}[1]{{\overrightarrow{#1}}}
\renewcommand{\d}{{\mathrm{d}}}
\providecommand{\degres}{{^\circ}}
\providecommand{\degre}{{^\circ}}
\providecommand{\blue}{\color{blue}}
\providecommand{\red}{\color{red}}
\providecommand{\enskip}{\quad}
\providecommand{\llbracket}{\lbrack\!\lbrack}
\providecommand{\rrbracket}{\rbrack\!\rbrack}
\providecommand{\euro}{{\text{€}}}
\providecommand{\no}{{\text{n°}}}
\providecommand{\interval}[2]{{[#1;#2]}}
\providecommand{\intervFO}[2]{{[#1;#2[}}
\providecommand{\dotfill}{\ldots}

\providecommand{\ds}{\displaystyle}
\providecommand{\textding}[1]{#1}
\providecommand{\Oij}{\(\mathrm{O};\vec{\imath},\vec{\jmath}\)}
\providecommand{\Oijk}{\(\mathrm{O};\vec{\imath},\vec{\jmath},\vec{k}\)}
\providecommand{\Ouv}{\(\mathrm{O};\vec{u},\vec{v}\)}
\providecommand{\intervFF}[2]{\ensuremath{[#1;#2]}}
\renewcommand{\euro}{\ifmmode\text{\texteuro}\else\texteuro\fi}
\DeclareUnicodeCharacter{00B0}{\ifmmode^{\circ}\else\textdegree\fi}
\DeclareUnicodeCharacter{20AC}{\ifmmode\text{\texteuro}\else\texteuro\fi}
% Énoncés : Éducation nationale. Transcription et corrigés : APMEP, Denis Vergès et contributeurs. Découpage et classement : CoopMaths, Sébastien Lozano, Éric Elter et contributeurs. Adaptation de la présentation : Maths Explorer.
% CC BY-SA : https://coopmaths.fr/www/about/
% https://coopmaths.fr/alea/static/eam/2025/tex/eam_specifique_2025_06_sujet0v1_2_cor.tex
\begin{document}
\noindent\textit{APMEP / CoopMaths. CC BY-SA. Présentation adaptée par Maths Explorer.}\par\medskip
\medskip

Albert a acquis un étang d'une surface de \np{2000}~m$^2$.
Le jour de son anniversaire, un dimanche, il installe des nénuphars sur une surface de $200$~m$^2$.



\begin{enumerate}
\item Le dimanche d'après, la surface des nénuphars a augmenté de $40$~m$^2$.
	\begin{enumerate}
		\item 
$\dfrac{40}{200}=\dfrac{20}{100}$ donc le pourcentage d'augmentation est de 20\,\%.		
				\item La surface occupée par les nénuphars est alors de 240~m$^2$.
	\end{enumerate}














\item Dans cette question, on suppose que la surface occupée par les nénuphars augmente de $40$ m$^2$ chaque semaine, depuis la date de l'anniversaire, tant que cela est possible.
	\begin{enumerate}
		\item $10\times 40=400$ et $200+400=600$
		
Donc la surface occupée par les nénuphars 10 semaines après l'anniversaire est de 600~m$^2$.
		\item 
Si  la surface occupée par les nénuphars est égale à $580$~m$^2$, cela veut dire qu'il y a eu une augmentation de 380~m$^2$ de la surface des nénuphars. Or 380 n'est pas un multiple de 40, donc ce n'est pas possible d'arriver à une surface occupée par les nénuphars de 580~m$^2$.
		
		\item
L'étang sera  entièrement recouvert de nénuphars quand la surface des nénuphars aura atteint \np{2000}~m$^2$, c'est-à-dire quand la surface des nénuphars aura augmenté de \np{1800}~m$^2$.

$\dfrac{\np{1800}}{40}=45$; il faudra donc 45 semaines pour que l'étang soit entièrement recouvert de nénuphars.
		
		
	\end{enumerate}
\item Dans cette question, on suppose que la surface occupée par les nénuphars augmente de $20\,\%$ chaque semaine, depuis la date de l'anniversaire, tant que cela est possible.
	\begin{enumerate}
		\item 
\begin{list}{\textbullet}{}
\item Une semaine après l'anniversaire la surface occupée par les nénuphars est, en m$^2$, de: $200+\dfrac{20}{100}\times 200 = 200+40=240$.
\item Deux semaines après l'anniversaire la surface occupée par les nénuphars est, en m$^2$, de: $240+\dfrac{20}{100}\times 240 = 240+48=288$.
\end{list}		
		
		\item On considère un entier naturel $n$. 
		
Ajouter 20\,\%, c'est multiplier par $1+\dfrac{20}{100}=1,2$.

Pour avoir la surface occupée par les nénuphars au bout d'une semaine, il faut multiplier la surface du départ par $1,2$.
		
Pour avoir la surface occupée par les nénuphars au bout de deux semaines, il faut multiplier la surface du départ par $1,2\times 1,2=1,2^2$. Etc.

Pour avoir la surface occupée par les nénuphars au bout de $n$ semaines, il faut multiplier la surface du départ par $1,2^n$.
		
La surface occupée par les nénuphars $n$ semaines après l'anniversaire est donc
$200\times 1,2^n$.
		
		\item 
L'étang sera entièrement recouvert par les nénuphars quand $n$ sera tel que $200\times 1,2^n \geqslant \np{2000}$, c'est-à-dire $1,2^n \geqslant 10$.

D'après le tableau fourni, ce sera pour $n=13$, donc au bout de 13 semaines.				
	\end{enumerate}
	
\begin{center}
\begin{tabularx}{\linewidth}{|l|*{9}{>{\centering \arraybackslash}X|}}\hline
$n=$& 0 &1 &2 &5 &10 &12 &13 &14 &15\\ \hline
$1,2^n \approx$& 1 &1,2& 1,44& 2,49& 6,19& 8,92& 10,70& 12,84& 15,40\\ \hline
\end{tabularx}
\end{center}

\item  Schéma sur lequel apparaissent l'allure des nuages de points traduisant la progression de la surface occupée par les nénuphars, aussi bien dans le cas de la question 2 (en noir) que dans le cas de la question 3 (en bleu).

\begin{center}
\scalebox{0.7}{
\psset{xunit=0.3cm,yunit=0.008cm,arrowsize=2pt 3}
\def\xmin{-5}   \def\xmax{47} \def\ymin{-100}   \def\ymax{2200}
\begin{pspicture*}(\xmin,\ymin)(\xmax,\ymax)
\psgrid[xunit=1.5cm,yunit=0.8cm,subgriddiv=1,  gridlabels=0, gridcolor=lightgray](0,0)(15,22)
\psaxes[arrowsize=3pt 3, ticksize=-2pt 2pt, Dx=5,Dy=100](0,0)(-4.9,-99)(\xmax,2199)
\psplot[plotpoints=46,plotstyle=dots]{0}{45}{200 40 x mul add}
\psline[linestyle=dashed](45,0)(45,2000)
\psplot[plotpoints=14,plotstyle=dots,linecolor=blue]{0}{13}{200 1.2 x exp mul}
\psline[linecolor=blue,linestyle=dashed](13,0)(13,2140)
\uput{8pt}[d](13,0){\blue 13}
\psplot[linecolor=red,linestyle=dashed]{0}{\xmax}{2000}
\uput[r](0,2150){Surface nénuphars} \uput[u](37,0){Nombre de semaines}
\end{pspicture*}
}
\end{center}

\end{enumerate}

\bigskip

\end{document}
